\documentclass[10pt]{mypackage}

%\usepackage{mlmodern}
%\usepackage{newpxtext,eulerpx,eucal}
%\renewcommand*{\mathbb}[1]{\varmathbb{#1}}

%\usepackage{homework}
%\usepackage{notes}

%\usepackage[ backend=bibtex, style = alphabetic, sorting=ynt ]{biblatex}
%\addbibresource{all_references.bib}

\usepackage{parskip}

\fancyhf{}
\fancyhead[R]{Avinash Iyer}
\fancyhead[L]{Hahn--Banach Theorems in Locally Convex Spaces}
\fancyfoot[C]{\thepage}

\setcounter{secnumdepth}{0}

\begin{document}
\RaggedRight
\section{Primer on Locally Convex Topological Vector Spaces}%
There are two primary notions we are concerned with when it comes to locally convex topological vector spaces: the notion of a topological vector space, and the notion of local convexity. We start by discussing topological vector spaces very quickly.
\begin{definition}
  Let $V$ be a vector space over $\C$. We say a topology $\tau$ on $V$ is compatible with the vector space structure if $\left( x,y \right) \mapsto x + y$ and $\left( \lambda,x \right)\mapsto \lambda x$ are continuous operations for all $x,y\in V$ and $\lambda\in \C$. If $V$ is equipped with a topology that is compatible with the vector space structure, then $V$ is called a topological vector space.
\end{definition}
\begin{definition}
  A subset $W\subseteq V$ is said to be \textit{absorbing} if for any $x\in V$, there is some $\lambda_0\in \C$ such that $x\in \lambda W$ for all $ \left\vert \lambda \right\vert\geq \left\vert \lambda_0 \right\vert $.

  We say the subset $C\subseteq V$ is \textit{circled} if $\lambda C \subseteq C$ for all $\left\vert \lambda \right\vert\leq 1$.

  We say the subset $A\subseteq V$ is \textit{convex} if, for any $x,y\in A$, we have $\lambda x + \left( 1-\lambda \right) y\in A$ for all $0 < \lambda < 1$.
\end{definition}
\begin{proposition}
  A topology $\tau$ on $V$ is compatible with the vector space structure if and only if $\tau$ is translation-invariant and possesses a neighborhood base $\mathcal{B}$ at $0$ satisfying
  \begin{enumerate}[(i)]
    \item for each $W\in \mathcal{B}$, there is $U\in \mathcal{B}$ such that $U + U \subseteq W$;
    \item each $W\in \mathcal{B}$ is circled and absorbing;
    \item there exists $0 < \left\vert \lambda \right\vert < 1$ such that whenever $W\in \mathcal{B}$, so too is $ \lambda W $.
  \end{enumerate}
\end{proposition}
\end{document}
